\begin{list}{}{\setlength{\leftmargin}{0.5in}\setlength{\rightmargin}{0.5in}}\item[]
Agentomics-ML runs machine-learning experiments as a state machine: the steps are fixed, each step's output is checked,
and the language model works only inside a step. A \emph{skill} describes the same procedure in a document that a
general-purpose agent reads and follows on its own. We ask whether the enforced structure still pays off, comparing the
two on the AGO2 CLASH miRNA--target classification task with nine language models, each run under both methods with
reasoning effort \texttt{none} and \texttt{high}: 36 runs, each planned for three iterations and scored on a held-out test set.

For top-of-line models with reasoning enabled the method makes no measurable difference: both score 0.902--0.910. Mid-range models
mostly do better with the state machine, which is clearly ahead in \PairsSm{} of the 18 model-and-setting pairs against
\PairsSkill{} for the skill, although the overall difference is not significant ($p=\WilcoxonP$). The state machine's
advantage is largest after a single iteration and shrinks as the agent iterates; two iterations give most of the final
score for about half the training time. The skill is cheaper and faster: \SkillCost{} and \SkillTime{} in total, against
\SmCost{} and \SmTime{}. Nothing forces an agent to run the skill's own test of each iteration, yet in sixteen of the
eighteen skill runs the agent ran it on every iteration and passed. Nine of the 36 runs were repeated with a second implementation; with
the original runs in their place the state machine is clearly ahead in \OrigPairsSm{} pairs and the skill in
\OrigPairsSkill{} ($p=\OrigWilcoxonP$).
\end{list}
